\chapter{LANDASAN TEORI}

\vspace{0.5cm}
\noindent\fbox{\begin{minipage}{\textwidth}
\textit{\textbf{Pedoman Penulisan:} Menguraikan teori-teori apa yang mendukung judul, berupa konsep dasar tentang teori tersebut yang berkaitan langsung dengan masalah yang diteliti. Ditulis juga tools/software komponen yang digunakan.}
\end{minipage}}
\vspace{0.5cm}

\section{Tinjauan Pustaka}

\vspace{0.5cm}
\noindent\fbox{\begin{minipage}{\textwidth}
\textit{\textbf{Pedoman Penulisan:} Berisi tentang semua teori-teori yang berhubungan dengan Skripsi yang akan dibahas.}
\end{minipage}}
\vspace{0.5cm}

\section{Penelitian Terkait}

\vspace{0.5cm}
\noindent\fbox{\begin{minipage}{\textwidth}
\textit{\textbf{Pedoman Penulisan:} Penjelasan tentang penelitian berdasarkan penelitian serupa dari jurnal-jurnal ilmiah. Minimal 5 kutipan jurnal 5 tahun terakhir (format APA).}
\end{minipage}}
\vspace{0.5cm}

\section{Tinjauan Organisasi / Objek Penelitian}
